\documentclass[11pt]{article}

\usepackage[letterpaper,margin=0.82in,headheight=15pt]{geometry}
\usepackage[T1]{fontenc}
\usepackage{lmodern}
\usepackage{microtype}
\usepackage{amsmath,amssymb,amsthm,bm,mathtools}
\usepackage{stmaryrd}
\usepackage{booktabs,longtable,tabularx,array,multirow}
\usepackage{enumitem}
\usepackage{xcolor}
\usepackage{fancyhdr}
\usepackage{graphicx}
\usepackage{tikz-cd}
\usepackage{hyperref}
\usepackage[nameinlink,noabbrev]{cleveref}

\definecolor{navy}{HTML}{17324D}
\definecolor{teal}{HTML}{147D92}
\definecolor{orange}{HTML}{D47A27}
\definecolor{softblue}{HTML}{EEF4F7}
\definecolor{softgray}{HTML}{F3F5F6}
\definecolor{darkgray}{HTML}{46545E}
\definecolor{softorange}{HTML}{FFF4E8}
\definecolor{softgreen}{HTML}{EDF7F0}
\definecolor{softred}{HTML}{FCEEEE}

\hypersetup{
  colorlinks=true,
  linkcolor=navy,
  citecolor=teal,
  urlcolor=teal,
  pdftitle={Volume XI: Microscopic Nonzero-Transfer GTMD Helicity Matrices, Projector Closure, and Wigner/OAM Consistency},
  pdfauthor={DeuteronWigner project}
}

\pagestyle{fancy}
\fancyhf{}
\fancyhead[L]{\small\color{darkgray}DeuteronWigner formalism - Volume XI}
\fancyhead[R]{\small\color{darkgray}Microscopic nonzero-transfer GTMD closure}
\fancyfoot[L]{\small\color{darkgray}31 July 2026}
\fancyfoot[R]{\small\color{darkgray}\thepage}
\renewcommand{\headrulewidth}{0.4pt}

\setlist[itemize]{leftmargin=1.5em,itemsep=2pt,topsep=3pt}
\setlist[enumerate]{leftmargin=1.7em,itemsep=2pt,topsep=3pt}
\setlength{\parindent}{0pt}
\setlength{\parskip}{5pt}
\setlength{\emergencystretch}{2em}
\renewcommand{\arraystretch}{1.16}
\allowdisplaybreaks

\newtheorem{definition}{Definition}[section]
\newtheorem{axiom}[definition]{Axiom}
\newtheorem{principle}[definition]{Design principle}
\newtheorem{requirement}[definition]{Requirement}
\newtheorem{proposition}[definition]{Proposition}
\newtheorem{theorem}[definition]{Theorem}
\newtheorem{lemma}[definition]{Lemma}
\newtheorem{corollary}[definition]{Corollary}
\newtheorem{remark}[definition]{Remark}
\newtheorem{decision}[definition]{Model decision}

\crefname{definition}{definition}{definitions}
\crefname{axiom}{axiom}{axioms}
\crefname{principle}{design principle}{design principles}
\crefname{requirement}{requirement}{requirements}
\crefname{proposition}{proposition}{propositions}
\crefname{theorem}{theorem}{theorems}
\crefname{lemma}{lemma}{lemmas}
\crefname{corollary}{corollary}{corollaries}
\crefname{remark}{remark}{remarks}
\crefname{decision}{model decision}{model decisions}

\newcommand{\kT}{\bm{k}_{T}}
\newcommand{\pT}{\bm{p}_{T}}
\newcommand{\qT}{\bm{q}_{T}}
\newcommand{\DT}{\bm{\Delta}_{T}}
\newcommand{\bD}{\bm{b}_{\Delta}}
\newcommand{\bM}{\bm{b}_{\mathrm{TMD}}}
\newcommand{\dd}{\mathrm{d}}
\newcommand{\Tr}{\operatorname{Tr}}
\newcommand{\tr}{\operatorname{tr}}
\newcommand{\diag}{\operatorname{diag}}
\newcommand{\rank}{\operatorname{rank}}
\newcommand{\Span}{\operatorname{span}}
\newcommand{\Ker}{\operatorname{Ker}}
\newcommand{\Ran}{\operatorname{Ran}}
\newcommand{\Cov}{\operatorname{Cov}}
\newcommand{\Var}{\operatorname{Var}}
\newcommand{\Id}{\operatorname{Id}}
\newcommand{\Hom}{\operatorname{Hom}}
\newcommand{\End}{\operatorname{End}}
\newcommand{\Hil}{\mathcal H}
\newcommand{\LF}{\ensuremath{\mathrm{LF}}}
\newcommand{\BLFQ}{\ensuremath{\mathrm{BLFQ}}}
\newcommand{\TMD}{\ensuremath{\mathrm{TMD}}}
\newcommand{\GTMD}{\ensuremath{\mathrm{GTMD}}}
\newcommand{\GPD}{\ensuremath{\mathrm{GPD}}}
\newcommand{\PDF}{\ensuremath{\mathrm{PDF}}}
\newcommand{\QCD}{\ensuremath{\mathrm{QCD}}}
\newcommand{\EFT}{\ensuremath{\mathrm{EFT}}}
\newcommand{\TTN}{\ensuremath{\mathrm{TTN}}}
\newcommand{\TN}{\ensuremath{\mathrm{TN}}}
\newcommand{\PCAC}{\ensuremath{\mathrm{PCAC}}}
\newcommand{\WTI}{\ensuremath{\mathrm{WTI}}}
\newcommand{\cO}{\mathcal O}
\newcommand{\cM}{\mathcal M}
\newcommand{\cR}{\mathcal R}
\newcommand{\cS}{\mathcal S}
\newcommand{\cT}{\mathcal T}
\newcommand{\cV}{\mathcal V}
\newcommand{\cW}{\mathcal W}
\newcommand{\cE}{\mathcal E}
\newcommand{\cG}{\mathcal G}
\newcommand{\cH}{\mathcal H}
\newcommand{\cK}{\mathcal K}
\newcommand{\cQ}{\mathcal Q}
\newcommand{\cP}{\mathcal P}
\newcommand{\cC}{\mathcal C}
\newcommand{\cZ}{\mathcal Z}
\newcommand{\cD}{\mathcal D}
\newcommand{\cN}{\mathcal N}
\newcommand{\cU}{\mathcal U}
\newcommand{\cF}{\mathcal F}
\newcommand{\cA}{\mathcal A}
\newcommand{\cB}{\mathcal B}
\newcommand{\cL}{\mathcal L}
\newcommand{\cJ}{\mathcal J}
\newcommand{\cI}{\mathcal I}
\newcommand{\eps}{\varepsilon}
\newcommand{\abs}[1]{\left|#1\right|}
\newcommand{\norm}[1]{\left\lVert#1\right\rVert}
\newcommand{\inner}[2]{\left\langle #1\middle|#2\right\rangle}
\newcommand{\avg}[1]{\left\langle #1\right\rangle}
\newcommand{\phys}{\mathrm{phys}}
\newcommand{\eff}{\mathrm{eff}}
\newcommand{\bare}{\mathrm{bare}}
\newcommand{\ind}{\mathrm{ind}}
\newcommand{\ct}{\mathrm{ct}}
\newcommand{\reg}{\mathrm{reg}}
\newcommand{\trunc}{\mathrm{trunc}}
\newcommand{\ren}{\mathrm{ren}}
\newcommand{\num}{\mathrm{num}}
\newcommand{\UV}{\mathrm{UV}}
\newcommand{\IR}{\mathrm{IR}}
\newcommand{\COM}{\mathrm{CM}}
\newcommand{\Lz}{L_z}
\newcommand{\Jz}{J^z}
\newcommand{\PV}{\operatorname{PV}}
\newcommand{\Hthree}{\Hil_{qqq}}
\newcommand{\Hfour}{\Hil_{qqqg}}
\newcommand{\Hfiveu}{\Hil_{qqqu\bar u}}
\newcommand{\Hfived}{\Hil_{qqqd\bar d}}
\newcommand{\AssumptionBundle}{\texttt{AssumptionBundle}}
\newcommand{\PredictionPlan}{\texttt{PredictionPlan}}
\newcommand{\StateBundle}{\texttt{H3MicroscopicStateBundle}}
\newcommand{\MomentumFiber}{\texttt{MomentumFiber}}
\newcommand{\RecoilMap}{\texttt{SymmetricXiZeroRecoil}}
\newcommand{\StateEvaluator}{\texttt{MomentumSpaceStateEvaluator}}
\newcommand{\OverlapKernel}{\texttt{MicroscopicGTMDOverlapKernel}}
\newcommand{\HelicityMatrix}{\texttt{PartonTargetHelicityMatrix}}
\newcommand{\ProjectorManifest}{\texttt{ProjectorClosureManifest}}
\newcommand{\MomentManifest}{\texttt{MomentClosureManifest}}
\newcommand{\WignerManifest}{\texttt{WignerOAMClosureManifest}}
\newcommand{\ParentBundle}{\texttt{MicroscopicGTMDParentBundle}}
\newcommand{\ReplacementManifest}{\texttt{AnalyticParentReplacementManifest}}
\newcommand{\ConvergenceManifest}{\texttt{H4ConvergenceManifest}}
\newcommand{\ProvComplex}{\texttt{Provenance2Complex}}
\newcommand{\Amp}{\mathbf{Amp}}
\newcommand{\Match}{\mathbf{Match}}
\newcommand{\Red}{\mathbf{Red}}
\newcommand{\Proc}{\mathbf{Proc}}
\newcommand{\Infer}{\mathbf{Infer}}

\newcolumntype{Y}{>{\raggedright\arraybackslash}X}
\newcolumntype{L}[1]{>{\raggedright\arraybackslash}p{#1}}

\title{\vspace{-1.0cm}
{\color{navy}\bfseries Volume XI: Microscopic Nonzero-Transfer GTMD Helicity Matrices,\\[2mm]
Projector Closure, and Wigner/OAM Consistency}\\[4mm]
\large Common H3 quark--antiquark--gluon parents, exact T-even reductions,
local-current and energy--momentum moments, and controlled replacement of
analytic overlap pilots}
\author{DeuteronWigner project}
\date{31 July 2026}

\begin{document}
\maketitle

\begin{center}
\begin{minipage}{0.94\textwidth}
\colorbox{softblue}{\parbox{0.96\textwidth}{
\textbf{Status and purpose.}
Volumes~0--X established the typed geometric architecture, regulated
light-front Hilbert spaces, common overlap maps, Wilson-line and cut pilots,
matched nuclear and QCD interfaces, shared inference, an assumption-conditioned
tensor-network compiler, and an interacting H3 nucleon state containing
\(qqq\), \(qqqg\), \(qqqu\bar u\), and \(qqqd\bar d\) sectors.  C11/H4 is
the active implementation package for the next transition: evaluation of
full nonzero-transfer quark, antiquark, and gluon helicity matrices from that
single microscopic state; complete leading-twist T-even projector closure;
commuting TMD, GPD, PDF, current, energy--momentum, Wigner, and OAM reductions;
and a scoped replacement of the earlier analytic common-parent pilots.  This
volume is the normative specification for H4.  It is intentionally restricted
to zero skewness and Wilson order zero.  It does not claim a soft-subtracted
QCD GTMD, a physical TMD, a complete continuum Hamiltonian trajectory, a
nuclear state, process factorization, or inference readiness.  A failure of a
projector or moment closure test must be repaired in the state evaluator,
operator basis, recoil map, current, or numerical implementation; it may not
be hidden by a GTMD-specific coefficient.
}}
\end{minipage}
\end{center}

\begin{abstract}
Generalized transverse-momentum dependent distributions are useful only when
they are retained as common off-forward operator matrix elements rather than
used as names for independently parameterized curves.  We formulate the
first microscopic DeuteronWigner H4 parent by evaluating the H3 Hamiltonian
eigenstate between distinct incoming and outgoing light-front momentum
fibers in the symmetric zero-skewness frame.  One versioned recoil map shifts
the active and spectator intrinsic transverse momenta, and one momentum-space
state evaluator converts exact BLFQ coefficients or a symmetry-adapted tree
tensor network into continuous amplitudes at the shifted coordinates.  The
same diagonal Fock-space overlap kernel generates positive-\(x\) quark,
antiquark, and gluon parents with common normalization, regulator, member,
phase, and assumption identity.

For quarks and antiquarks, the full target--parton helicity matrix is retained
and projected onto the complete leading-twist spin-\(\tfrac12\) GTMD basis
\(F_{1,1\ldots4}\), \(G_{1,1\ldots4}\), and
\(H_{1,1\ldots8}\).  Projectors are constructed from the Gram matrix of the
declared covariant basis.  Full-rank extraction is required only at generic
nondegenerate kinematics; at the forward or collinear boundaries the
identifiable reduced basis is reported instead of using a pseudoinverse to
invent separated amplitudes.  The direct positive-\(x\) antiquark matrices
retain operator-specific charge-conjugation signatures.  The gluon parent
retains its complete transverse tensor, target and gluon helicities, both
microscopic \(qqqg\) color outer multiplicities, and exact trace,
helicity-antisymmetric, and symmetric-traceless reconstruction.  At Wilson
order zero every link-odd projection is an exact null test.

TMD, GPD, PDF, local-current, energy--momentum, Wigner, and OAM quantities are
defined as typed reductions of the same parent.  Vector, axial, tensor, and
gluon moments use their correct quark--antiquark signs and gluon normalization.
The phase-space OAM moment is checked against a transfer derivative and the
direct H3 canonical operator, with all Fourier and gauge-path conventions
retained.  Forward Gram positivity is distinguished from nonforward
overlap bounds and from Wigner quasiprobability.  Exact, Krylov, and TTN
state representations, transfer grids, quadratures, basis and Fock towers,
and bond dimensions enter separate convergence ledgers.  Finally, the H4
parent replaces the C3/C4 analytic parents only inside a typed validation
scope after all closure gates pass; the analytic models remain immutable
oracles and are not fitting targets.  The package specifies formal
requirements, software objects, schemas, benchmark families, negative tests,
provenance cells, and acceptance criteria for the first microscopic
nonzero-transfer common parent.
\end{abstract}

\tableofcontents

\section{Scientific problem, scope, and interfaces}
\label{sec:scope}

\subsection{The H4 transition}

H3 supplies one normalized interacting state
\begin{equation}
 |N,\Lambda;s,r,\mathfrak A\rangle
 =
 |\Psi_{qqq}^{\Lambda}\rangle
 +|\Psi_{qqqg}^{\Lambda}\rangle
 +|\Psi_{qqqu\bar u}^{\Lambda}\rangle
 +|\Psi_{qqqd\bar d}^{\Lambda}\rangle,
 \label{eq:h3-state}
\end{equation}
where \(s\) is the microscopic member, \(r\) is the regulator and resolution
identity, and \(\mathfrak A\) is the assumption bundle.  H3 already closes
state probability, valence flavor, charge, baryon number, longitudinal
momentum, and a finite-basis canonical \(\Jz\) ledger.  It also supplies
Hamiltonian-consistent vector, axial, pseudoscalar, field-strength, and OAM
operators.  What it does not yet provide is one complete off-forward partonic
parent evaluated from the actual state.

The H4 transition is
\begin{equation}
\begin{aligned}
 &\bigl\{|N\rangle,\cO_q,\cO_{\bar q},\cO_g,\cJ,\cT^{\mu\nu},\cL_z\bigr\}
 \\
 &\qquad\longrightarrow
 \bigl\{\cW_q,\cW_{\bar q},\cW_g\bigr\}(x,\kT,\DT)
 \\
 &\qquad\longrightarrow
 \bigl\{\TMD,\GPD,\mathrm{PDF},\mathrm{current},\mathrm{EMT},
 \mathrm{Wigner},\mathrm{OAM}\bigr\}.
\end{aligned}
\label{eq:h4-chain}
\end{equation}

\begin{decision}[H4 scientific boundary]
H4 is a regulated zero-skewness, Wilson-order-zero common-parent
calculation.  It may issue a microscopic nonzero-transfer closure status.  It
may not issue physical QCD GTMD, TMD, process, nuclear, evolution, or
inference readiness.
\end{decision}

\subsection{Inputs inherited without reinterpretation}

The following objects are normative inputs:
\begin{enumerate}
\item the H3 Hamiltonian, renormalization trajectory, state bundles, exact and
TTN state representations, and phase conventions;
\item the C1--C2 typed coordinate, rank, Wilson-path, operator-identity,
reduction, and provenance interfaces;
\item the C3 symmetric zero-skewness recoil authority and analytic overlap
oracles;
\item the C4 regulated reduction and current/EMT route semantics;
\item the C5--C6 Wilson, cut, ordered-link, color, and phase-budget types,
which remain inactive at Wilson order zero;
\item the Volume~V matching and evolution gates and the Volume~VI inference
gates, which remain fail-closed.
\end{enumerate}

H4 may extend these interfaces but may not create a parallel coordinate,
recoil, operator, rank, path, or provenance system.

\subsection{Explicit nonclaims}

The following are not established by H4:
\begin{itemize}
\item continuum regulator independence;
\item nonzero skewness or ERBL overlap structure;
\item polynomiality beyond the tested zero-skewness moments;
\item a fully soft-subtracted and rapidity-renormalized GTMD;
\item physical link-odd Sivers, Boer--Mulders, or gluonic-pole functions;
\item full non-Abelian Ward or Slavnov--Taylor closure;
\item a microscopic deuteron or spin-1 target state;
\item physical TMD evolution, process factorization, or a global posterior.
\end{itemize}

\section{Zero-skewness off-forward kinematics}
\label{sec:kinematics}

\subsection{Symmetric frame and momentum fibers}

Let
\begin{equation}
 p=P-\frac{\Delta}{2},
 \qquad
 p'=P+\frac{\Delta}{2},
 \qquad
 \Delta^+=0,
 \qquad
 \xi=-\frac{\Delta^+}{2P^+}=0,
 \label{eq:symmetric-frame}
\end{equation}
with
\begin{equation}
 t=\Delta^2=-\DT^2
\end{equation}
in the purely transverse frame.  The initial and final hadron spaces are
separate fibers,
\begin{equation}
 \Hil_{P,\Lambda}^{(r)}
 \qquad\text{and}\qquad
 \Hil_{P',\Lambda'}^{(r)}.
\end{equation}
An off-forward operator is therefore typed as
\begin{equation}
 \cO_{\Gamma,\gamma}^{a}(P',P):
 \Hil_{P}^{(r)}\longrightarrow\Hil_{P'}^{(r)},
 \label{eq:operator-fiber-map}
\end{equation}
not as an endomorphism of one nominal momentum space.

\begin{definition}[Momentum-fiber identity]
A \MomentumFiber{} stores the hadron momentum, light-front convention,
normalization, regulator, helicity basis, phase convention, and comparison-map
identity.  Two fibers are composable only through an explicit transfer or
matching map.
\end{definition}

\subsection{Single authoritative recoil map}

For an \(n\)-parton state and active parton \(j\), let
\(x_i>0\), \(\sum_i x_i=1\), and \(\sum_i\bm k_{iT}=0\).  The symmetric
zero-skewness recoil map is
\begin{align}
 \bm\kappa_{jT}^{\mathrm{in}}
 &=\bm k_{jT}-\frac{1-x_j}{2}\DT,
 &
 \bm\kappa_{jT}^{\mathrm{out}}
 &=\bm k_{jT}+\frac{1-x_j}{2}\DT,
 \label{eq:active-recoil}\\
 \bm\kappa_{iT}^{\mathrm{in}}
 &=\bm k_{iT}+\frac{x_i}{2}\DT,
 &
 \bm\kappa_{iT}^{\mathrm{out}}
 &=\bm k_{iT}-\frac{x_i}{2}\DT,
 \qquad i\neq j.
 \label{eq:spectator-recoil}
\end{align}

\begin{proposition}[Intrinsic closure]
The map in \cref{eq:active-recoil,eq:spectator-recoil} satisfies
\begin{equation}
 \sum_i\bm\kappa_{iT}^{\mathrm{in}}=0,
 \qquad
 \sum_i\bm\kappa_{iT}^{\mathrm{out}}=0,
\end{equation}
and has unit Jacobian in the independent intrinsic coordinates.
\end{proposition}

\begin{proof}
For the incoming state, the transfer-dependent contribution is
\begin{equation}
 -\frac{1-x_j}{2}\DT
 +\sum_{i\neq j}\frac{x_i}{2}\DT
 =0,
\end{equation}
since \(\sum_{i\neq j}x_i=1-x_j\).  The outgoing relation follows with the
opposite sign.  The map is affine with identity linear part on the independent
intrinsic coordinates, so its Jacobian is one.
\end{proof}

\begin{proposition}[Physical transfer and spectator preservation]
When intrinsic momenta are combined with the hadron transverse momenta
\(\bm p_T=-\DT/2\) and \(\bm p_T'=+\DT/2\), the active parton receives the
full physical transfer \(\DT\), while every spectator physical momentum is
unchanged.
\end{proposition}

\begin{requirement}[No local recoil formulas]
The \RecoilMap{} is the only implementation authority for
\cref{eq:active-recoil,eq:spectator-recoil}.  State evaluators, projectors,
species selectors, current reductions, and Wigner transforms may consume the
map but may not reproduce its factors locally.
\end{requirement}

\subsection{Transfer reversal and forward limit}

The transfer-reversal map is
\begin{equation}
 \cR_\Delta:
 (P',P,\DT,\Lambda',\Lambda)
 \longmapsto
 (P,P',-\DT,\Lambda,\Lambda').
\end{equation}
It exchanges incoming and outgoing recoil coordinates exactly.  At
\(\DT=0\), all shifted coordinates reduce to the common forward intrinsic
configuration.  These are exact property tests, not floating-point
approximations.

\section{Momentum-space realization of the microscopic state}
\label{sec:state-evaluator}

\subsection{Basis expansion}

For sector \(\nu\), write the H3 state as
\begin{equation}
 |\Psi_{\nu,\Lambda}^{(r,s)}\rangle
 =
 \sum_{\alpha\in\cB_{\nu,r}}
 C_{\nu\alpha}^{\Lambda;(r,s)}
 |\nu,\alpha;r\rangle,
 \label{eq:basis-state}
\end{equation}
where \(\alpha\) contains longitudinal modes, transverse radial and orbital
indices, parton helicities, flavor, color outer multiplicity, permutation
representation, and center-of-mass identity.  The momentum-space amplitude is
\begin{equation}
 \psi_{\nu,\Lambda}^{(r,s)}
 \bigl(\{x_i,\bm\kappa_{iT},\lambda_i,c_i,f_i\}\bigr)
 =
 \sum_{\alpha}
 C_{\nu\alpha}^{\Lambda;(r,s)}
 \Phi_{\nu\alpha}^{(r)}
 \bigl(\{x_i,\bm\kappa_{iT},\lambda_i,c_i,f_i\}\bigr).
 \label{eq:momentum-state}
\end{equation}

The basis function \(\Phi_{\nu\alpha}^{(r)}\) includes the declared
longitudinal normalization, two-dimensional oscillator or alternate
transverse basis, exact color tensor, signed fermion permutation tensor, and
center-of-mass factor.  The oscillator scale is a basis coordinate, not a
TMD width.

\begin{definition}[Momentum-space state evaluator]
The \StateEvaluator{} evaluates \cref{eq:momentum-state} from either the exact
basis coefficient vector or the symmetry-adapted TTN.  Its result identity
contains the state-bundle hash, representation type, bond dimension, basis
resolution, phase convention, coordinate fiber, and numerical tolerance.
\end{definition}

\subsection{Exact and TTN representations}

The full-bond TTN and the exact basis vector represent the same finite state.
For every supported configuration \(z\), H4 requires
\begin{equation}
 \psi_{\nu}^{\mathrm{TTN},\chi_{\mathrm{full}}}(z)
 =
 \psi_{\nu}^{\mathrm{exact}}(z)
 +\delta_{\mathrm{repr}}(z),
 \qquad
 \norm{\delta_{\mathrm{repr}}}\leq\eps_{\mathrm{repr}}.
 \label{eq:exact-ttn-amplitude}
\end{equation}
Finite-bond results retain a separate truncation envelope and cannot borrow
the full-bond identity.

\begin{requirement}[Observable-sensitive TTN convergence]
Bond convergence shall be reported for matrix elements and interference
blocks, not only for the energy.  A bond space that reproduces the mass but
removes the only \(\Delta L_z\) block supporting a requested GTMD is not
admissible for that GTMD.
\end{requirement}

\subsection{Quadrature and continuous coordinates}

The finite basis defines a continuous momentum-space amplitude, but the
numerical matrix element still requires integration.  A quadrature manifest
shall contain:
\begin{itemize}
\item longitudinal simplex rule and endpoint policy;
\item transverse radial and azimuthal grids;
\item transfer grid and angular orientations;
\item interpolation or basis-evaluation order;
\item independent high-order or analytic oracle where available;
\item absolute and relative errors for each reduction route.
\end{itemize}
The quadrature resolution is a numerical axis distinct from the Hamiltonian
basis resolution.

\subsection{Derivatives with respect to transfer}

OAM and impact-parameter moments require derivatives near \(\DT=0\).  H4 may
use analytic basis derivatives, automatic differentiation, complex-step
methods where compatible, or symmetric finite differences.  The derivative
method and step-size study are part of the result identity.  A finite
difference that mixes states with inconsistent phase conventions is invalid.

\section{Common microscopic overlap kernel}
\label{sec:overlap}

\subsection{Zeroth-rescattering diagonal overlap}

At \(\xi=0\), positive \(x\), and Wilson order zero, the canonical leading-
twist bilocal operator produces a diagonal retained-parton-number overlap in
the DGLAP region.  For species \(a\),
\begin{align}
 W_{a/N,\Lambda'\Lambda}^{[\Gamma],(r,s)}
 (x,\kT,\DT)
 ={}&
 \sum_{\nu\in\cF_{\mathrm{H3}}}
 \sum_{j\in a,\nu}
 \int[\dd x]_{\nu}[\dd^2\bm k_T]_{\nu}
 \;\delta(x-x_j)
 \delta^{(2)}(\kT-\bm k_{jT})
 \nonumber\\
 &\times
 \psi_{\nu,\Lambda'}^{*(r,s)}
 \bigl(\{x_i,\bm\kappa_{iT}^{\mathrm{out}}\}\bigr)
 \cP_{j}^{[\Gamma]}
 \psi_{\nu,\Lambda}^{(r,s)}
 \bigl(\{x_i,\bm\kappa_{iT}^{\mathrm{in}}\}\bigr),
 \label{eq:common-overlap}
\end{align}
with exact spectator helicity, flavor, and color matching.

The sectors are
\begin{equation}
 \cF_{\mathrm{H3}}
 =
 \{qqq,qqqg,qqqu\bar u,qqqd\bar d\}.
\end{equation}
Quarks receive active slots from all sectors that contain the requested
flavor; antiquarks receive active positive-\(x\) slots only from the pair
sectors; gluons receive active slots only from \(qqqg\).

\begin{definition}[Overlap-kernel identity]
The \OverlapKernel{} contains the operator identity, incoming and outgoing
fibers, active-slot selector, recoil map, state member, regulator, Fock
sector, color multiplicity, helicity kernel, permutation convention,
quadrature, Wilson order, and matching status.
\end{definition}

\subsection{Number-changing and zero-mode terms}

H4 does not silently include off-diagonal Fock overlaps.  A term with
\(\nu'\neq\nu\) requires a named source such as:
\begin{itemize}
\item nonzero skewness and ERBL support;
\item an explicit Wilson-line expansion;
\item a constrained-field or zero-mode operator;
\item an induced effective operator generated by sector elimination;
\item operator renormalization mixing.
\end{itemize}
All such terms are outside the central H4 overlap unless explicitly declared
as a separate benchmark.  Missing support is reported as
\texttt{UNAVAILABLE\_AT\_THIS\_TRUNCATION}, never as a physical zero.

\subsection{Normalization and common-member propagation}

One microscopic member generates all species and reductions.  For member
\(s\),
\begin{equation}
 s
 \longmapsto
 \left\{
 W_{u,s},W_{d,s},W_{\bar u,s},W_{\bar d,s},W_{g,s}
 \right\}
 \longmapsto
 \left\{
 F_{A,s},\cO_{i,s}
 \right\}.
 \label{eq:member-propagation}
\end{equation}
The quark, antiquark, and gluon normalizations are therefore linked by state
probability, valence number, longitudinal momentum, helicity, and angular-
momentum ledgers.  They cannot be normalized independently after overlap
evaluation.

\section{Joint target--parton helicity matrices}
\label{sec:helicity-matrices}

The scalar coefficients conventionally called GTMDs are coordinates on a
larger helicity object.  H4 therefore stores the matrix before any named
projection is performed.  This reverses the common but unsafe workflow in
which independently generated scalar functions are later assembled into a
matrix and repaired if a symmetry or positivity condition fails.

\subsection{Quark and antiquark matrices}

For a quark or antiquark species, define the ordered joint index
\begin{equation}
 \alpha=(\lambda_a,\Lambda),
 \qquad
 \lambda_a\in\left\{+\tfrac12,-\tfrac12\right\},
 \qquad
 \Lambda\in\left\{+\tfrac12,-\tfrac12\right\}.
\end{equation}
The microscopic quark matrix is
\begin{equation}
 \cM^{q}_{\alpha'\alpha}
 (x,\kT,\DT)
 =
 \cM^{q}_{\lambda_q'\Lambda';\lambda_q\Lambda}
 (x,\kT,\DT),
 \label{eq:quark-helicity-matrix}
\end{equation}
and the antiquark matrix
\begin{equation}
 \cM^{\bar q}_{\alpha'\alpha}
 (x,\kT,\DT)
 =
 \cM^{\bar q}_{\lambda_{\bar q}'\Lambda';
                  \lambda_{\bar q}\Lambda}
 (x,\kT,\DT)
 \label{eq:antiquark-helicity-matrix}
\end{equation}
uses direct positive-\(x\) antiquark slots.  The canonical basis ordering is
stored in the object identity and is never inferred from matrix shape.

The entries of \cref{eq:quark-helicity-matrix} are obtained by replacing the
active kernel in \cref{eq:common-overlap} with the elementary quark-helicity
matrix unit
\begin{equation}
 \cK^{\lambda_q'\lambda_q}_{j}
 =|\lambda_q'\rangle\langle\lambda_q|
 \otimes\Id_{\rm spectator}.
\end{equation}
The standard Dirac projections are then fixed linear combinations:
\begin{align}
 \cM^{[\gamma^+]}&=\cM^{++}+\cM^{--},\nonumber\\
 \cM^{[\gamma^+\gamma_5]}&=\cM^{++}-\cM^{--},\nonumber\\
 \cM^{[i\sigma^{R+}\gamma_5]}&=2\cM^{+-},\qquad
 \cM^{[i\sigma^{L+}\gamma_5]}=2\cM^{-+},
 \label{eq:quark-helicity-dirac-map}
\end{align}
up to the declared spinor and circular-basis convention.  The adapter carrying
those factors is versioned and tested against explicit spinor bilinears.

\begin{requirement}[Complete matrix export]
\label{req:complete-quark-matrix}
Every supported quark and antiquark result must export all sixteen complex
entries of the \(4\times4\) joint matrix, including entries that vanish in the
selected truncation.  A vanishing entry records its exact selection-rule or
amplitude-content reason.
\end{requirement}

\subsection{Gluon matrices and transverse tensor}

For an active gluon, one may use circular helicities
\(\lambda_g=\pm1\) or transverse Cartesian indices \(i,j=1,2\).  H4 stores both
representations through one invertible adapter:
\begin{equation}
 \cM^{g}_{\lambda_g'\Lambda';\lambda_g\Lambda}
 \longleftrightarrow
 W^{g,ij}_{\Lambda'\Lambda}.
 \label{eq:gluon-circular-cartesian}
\end{equation}
The complete tensor has target-helicity shape \(2\times2\) and active-gluon
transverse shape \(2\times2\).  It is decomposed only after evaluation:
\begin{align}
 W_U&=\delta_{T,ij}W^{ij},\label{eq:gluon-trace}\\
 W_H&=i\epsilon_{T,ij}W^{ij},\label{eq:gluon-antisym}\\
 W_L^{ij}&=W^{(ij)}-\frac12\delta_T^{ij}W^k_{\ k}.
 \label{eq:gluon-st}
\end{align}
The reconstruction identity is
\begin{equation}
 W^{ij}
 =\frac12\delta_T^{ij}W_U
 -\frac{i}{2}\epsilon_T^{ij}W_H
 +W_L^{ij}.
 \label{eq:gluon-polarization-reconstruction}
\end{equation}

The two independent \(qqqg\) color-singlet outer multiplicities inherited from
H2/H3 remain separate state labels.  They are not the \(f^{abc}\)- and
\(d^{abc}\)-type gluonic-pole classes of a rescattering operator.  The former
classify the hadron state coupling; the latter classify a future Wilson-line
color contraction.  Conflating them is a type error.

\begin{requirement}[One gluon tensor parent]
\label{req:one-gluon-tensor-parent}
Trace, helicity-antisymmetric, and symmetric-traceless gluon projections must
reconstruct one common \(W^{ij}\).  Separate numerical parents for the three
polarization sectors are inadmissible.
\end{requirement}

\subsection{Target-helicity phase convention}

For both quarks and gluons, the target basis is ordered
\begin{equation}
 \{|+\tfrac12\rangle,|-\tfrac12\rangle\}.
\end{equation}
The state phase at each resolution is inherited from H3 state tracking: the
largest common nonzero component is real and positive after transport from the
coarser state.  At a degeneracy, the entire tracked subspace is parallel
transported and the matrix is transformed covariantly.  Individual eigenvector
phases are not independently reset at every transfer point.

\begin{proposition}[Basis covariance]
Let \(U_T\) and \(U_a\) be unitary changes of target and active-parton helicity
basis.  Then
\begin{equation}
 \cM\mapsto
 (U_a\otimes U_T)\cM(U_a\otimes U_T)^\dagger
 \label{eq:matrix-basis-covariance}
\end{equation}
leaves every scalar extracted with the correspondingly transformed projector
invariant.
\end{proposition}

\section{Complete quark and antiquark GTMD projector system}
\label{sec:quark-projectors}

\subsection{Generic leading-twist basis}

At zero skewness and nonzero transverse transfer, the complete leading-twist
quark basis for a spin-\(\tfrac12\) target contains
\begin{equation}
 \{F_{1,n}\}_{n=1}^{4},\qquad
 \{G_{1,n}\}_{n=1}^{4},\qquad
 \{H_{1,n}\}_{n=1}^{8},
 \label{eq:quark-16-basis}
\end{equation}
within a declared Lorentz, spinor, and transverse-tensor convention
\cite{Meissner2009}.  H4 treats these names as adapter coordinates on the
microscopic matrix, not as sixteen independent model functions.

A convenient reference parameterization is
\begin{equation}
 W^{[\gamma^+]}
 =\frac{1}{2M}\bar u(p',\Lambda')\,\cF\,u(p,\Lambda),
\end{equation}
with
\begin{equation}
 \cF=
 F_{1,1}
 +\frac{i\sigma^{i+}k_T^i}{P^+}F_{1,2}
 +\frac{i\sigma^{i+}\Delta_T^i}{P^+}F_{1,3}
 +\frac{i\sigma^{ij}k_T^i\Delta_T^j}{M^2}F_{1,4}.
 \label{eq:h4-F-basis}
\end{equation}
The axial sector is
\begin{equation}
 W^{[\gamma^+\gamma_5]}
 =\frac{1}{2M}\bar u(p',\Lambda')\,\cG\,u(p,\Lambda),
\end{equation}
with
\begin{equation}
 \cG=
 -\frac{i\epsilon_T^{ij}k_T^i\Delta_T^j}{M^2}G_{1,1}
 +\frac{i\sigma^{i+}\gamma_5k_T^i}{P^+}G_{1,2}
 +\frac{i\sigma^{i+}\gamma_5\Delta_T^i}{P^+}G_{1,3}
 +i\sigma^{+-}\gamma_5G_{1,4}.
 \label{eq:h4-G-basis}
\end{equation}
For the transverse tensor sector,
\begin{equation}
 W^{[i\sigma^{j+}\gamma_5]}
 =\frac{1}{2M}\bar u(p',\Lambda')\,\cH^j\,u(p,\Lambda),
\end{equation}
where
\begin{align}
 \cH^j={}&
 -\frac{i\epsilon_T^{ij}k_T^i}{M}H_{1,1}
 -\frac{i\epsilon_T^{ij}\Delta_T^i}{M}H_{1,2}
 +\frac{Mi\sigma^{j+}\gamma_5}{P^+}H_{1,3}
 +\frac{k_T^ji\sigma^{i+}\gamma_5k_T^i}{MP^+}H_{1,4}
 \nonumber\\
 &+\frac{\Delta_T^ji\sigma^{i+}\gamma_5k_T^i}{MP^+}H_{1,5}
 +\frac{\Delta_T^ji\sigma^{i+}\gamma_5\Delta_T^i}{MP^+}H_{1,6}
 +\frac{k_T^ji\sigma^{+-}\gamma_5}{M}H_{1,7}
 +\frac{\Delta_T^ji\sigma^{+-}\gamma_5}{M}H_{1,8}.
 \label{eq:h4-H-basis}
\end{align}
All factors and phases belong to a versioned basis adapter and are verified by
helicity-matrix reconstruction.

For each Dirac sector \(J\), construct a set of basis matrices
\begin{equation}
 \cB_J(\kT,\DT)
 =\{B_{J,A}(\kT,\DT)\}_{A=1}^{N_J},
 \end{equation}
where \(N_{\gamma^+}=4\), \(N_{\gamma^+\gamma_5}=4\), and
\(N_{i\sigma^{i+}\gamma_5}=8\).  The matrices include the declared factors of
\(i\), nucleon mass, \(k_R=k_x+ik_y\), \(\Delta_R=\Delta_x+i\Delta_y\), and
the spinor normalization.  These factors are not reconstructed from a paper
name at runtime.

The Gram matrix and dual projectors are
\begin{align}
 G_{AB}^{(J)}(z)
 &=\Tr\left[B_{J,A}^\dagger(z)B_{J,B}(z)\right],
 \label{eq:gtmd-gram}\\
 P_{J,A}(z)
 &=\sum_B\left[G^{(J)}(z)^{-1}\right]_{AB}B_{J,B}^\dagger(z),
 \label{eq:gtmd-dual-projector}
\end{align}
with \(z=(x,\kT,\DT)\) and inversion restricted to the identifiable subspace.
The microscopic coefficient is
\begin{equation}
 X_{J,A}(z)
 =\Tr\left[P_{J,A}(z)\cM_J(z)\right].
 \label{eq:gtmd-project}
\end{equation}

\begin{theorem}[Projector reconstruction]
\label{thm:projector-reconstruction}
At a kinematic point where \(G^{(J)}\) has full rank,
\begin{equation}
 \cM_J(z)
 =\sum_A X_{J,A}(z)B_{J,A}(z)
 \label{eq:projector-reconstruct}
\end{equation}
if and only if the matrix lies in the declared leading-twist basis span.
The reconstruction residual is a direct diagnostic of missing basis content,
a convention mismatch, or a higher-twist/numerical contribution.
\end{theorem}

\subsection{Conditioning and degeneracy}

At \(\DT=0\), \(\kT=0\), collinear configurations, or special relative
azimuths, the generic sixteen-coordinate basis becomes rank deficient.  H4
records
\begin{equation}
 \rank G^{(J)}(z),\qquad
 \kappa(G^{(J)}(z)),\qquad
 \sigma_{\min}^{\rm retained}(z),
 \label{eq:gram-diagnostics}
\end{equation}
and switches to a predefined reduced basis when the smallest singular value
falls below the declared conditioning threshold.

\begin{requirement}[No pseudoinverse overclaim]
\label{req:no-pseudoinverse-overclaim}
A Moore--Penrose pseudoinverse may be used to reconstruct the identifiable
matrix component, but its null-space coordinates must be reported as
\texttt{UNIDENTIFIABLE\_AT\_THIS\_KINEMATIC\_POINT}.  It may not be used to
publish all sixteen scalar GTMDs as if they were separately determined.
\end{requirement}

Kinematic rank loss is distinct from a dynamical zero.  The former means two
basis tensors become indistinguishable; the latter means the state supplies no
amplitude for an identifiable projector.

\subsection{Antiquark projector and charge-conjugation adapter}

The direct positive-\(x\) antiquark matrix is projected with the same geometric
basis but with an explicit charge-conjugation adapter
\begin{equation}
 \mathfrak C_J=
 \left(
 J,\eta_J,\text{endpoint order},\text{spinor phase},
 \bar3\text{ color action}
 \right),
 \label{eq:charge-conjugation-adapter}
\end{equation}
where \(\eta_J\) is the local-operator charge-conjugation signature.  This
adapter controls the relative quark/antiquark signs in local moments.  It does
not identify \(W_{\bar q}(x)\) with \(W_q(-x)\) inside the computational state.

\subsection{Standard-basis adapters}

The internal basis is versioned independently of literature notation.  A
standard-basis adapter supplies a matrix
\begin{equation}
 \bm X_{\rm reference}
 =\bm A_{\rm reference\leftarrow internal}\bm X_{\rm internal}
 \label{eq:reference-adapter}
\end{equation}
with an inverse on the full-rank generic space.  Every adapter declares:
\begin{itemize}
\item light-front and spinor normalization;
\item definitions of \(k_{R,L}\) and \(\Delta_{R,L}\);
\item metric and epsilon conventions;
\item mass factors and dimensions;
\item Hermiticity and time-reversal phases;
\item the reference equation and version.
\end{itemize}
An adapter is accepted only after random-coefficient round-trip tests and
explicit helicity-amplitude comparisons.

\section{Complete gluon tensor and helicity projector system}
\label{sec:gluon-projectors}

\subsection{Microscopic gluon operator space}

At Wilson order zero, the active-gluon H4 parent is a color-singlet hadron
matrix element of the linear field-strength operator supported by the H3
\(qqqg\) sector.  It is not yet the complete nonlinear continuum field-strength
operator.  The declared basis therefore records the field-strength
completeness status
\begin{equation}
 \texttt{LINEAR\_ACTIVE\_GLUON\_H4}
 \end{equation}
and preserves unresolved nonlinear, constrained-field, and higher-Fock
components as nonzero obligations.

Let \(C_{g,A}^{ij;\Lambda'\Lambda}(\kT,\DT)\) denote the allowed tensor basis
at the selected twist, parity, and Wilson order.  Define
\begin{equation}
 G^{g}_{AB}
 =\sum_{i,j,\Lambda',\Lambda}
 C_{g,A}^{ij;\Lambda'\Lambda *}
 C_{g,B}^{ij;\Lambda'\Lambda}.
 \label{eq:gluon-gram}
\end{equation}
The projectors follow from the inverse on the identifiable subspace, exactly as
in \cref{eq:gtmd-dual-projector}.

\begin{requirement}[Measured basis rank]
\label{req:measured-gluon-rank}
The number of independently extractable gluon coefficients is the numerical
and symbolic rank of \cref{eq:gluon-gram} under the declared symmetries and
kinematics.  It is not inferred from the number of names in an external table
or from the shape of an output array.
\end{requirement}

\subsection{Polarization-first closure}

The tensor is first decomposed by \cref{eq:gluon-trace,eq:gluon-antisym,eq:gluon-st}, and each component is then resolved in target helicity and
transverse harmonics.  Equivalently, one may use the circular-helicity basis
and transform afterward.  Both routes must agree:
\begin{equation}
 \Red_{\rm target}\circ\Red_{\rm gluon\ pol}[W]
 =
 \Red_{\rm gluon\ pol}\circ\Red_{\rm target}[W].
 \label{eq:gluon-projector-commute}
\end{equation}
Failure of this finite-dimensional square signals an index-order, convention,
or incomplete-basis defect.

\subsection{State color multiplicity versus Wilson color class}

The H4 state color label
\begin{equation}
 \alpha_{8}\in\{\rho,\lambda\}
\end{equation}
identifies the two ways the \(qqq\) octet couples with the active adjoint gluon
to a hadron singlet.  A future Wilson calculation introduces a distinct
three-adjoint operator contraction
\begin{equation}
 \cC\in\{f,d\}.
\end{equation}
The total future identity is therefore the ordered pair
\begin{equation}
 (\alpha_8,\cC),
 \label{eq:color-two-layer-id}
\end{equation}
not a single two-valued label.  H4 has \(\alpha_8\) but no active \(\cC\),
because Wilson order is zero.

\begin{requirement}[No premature gluonic-pole label]
\label{req:no-premature-fd}
An H4 gluon GTMD may not be labeled \(f\)-type or \(d\)-type.  Those labels
become defined only after an ordered two-link rescattering operator is attached
in H5.
\end{requirement}

\section{Discrete symmetries and exact H4 null sectors}
\label{sec:discrete-symmetry}

\subsection{Hermiticity}

For a Hermitian bilocal operator with the declared endpoint ordering,
\begin{equation}
 W_{\Lambda'\Lambda}^{a[J]}
 (x,\kT,\DT)
 =
 \left[
 W_{\Lambda\Lambda'}^{a[J]}
 (x,\kT,-\DT)
 \right]^*
 \label{eq:h4-hermiticity}
\end{equation}
up to the operator-specific index transpose.  The test is applied to the full
matrix before scalar projection and then independently to the extracted
coefficients.

\subsection{Light-front parity}

A typed parity adapter transforms target and parton helicities, transverse
momenta, spinor phases, and tensor indices together.  In schematic circular
notation,
\begin{equation}
 (\lambda_a,\Lambda,k_R,\Delta_R)
 \longmapsto
 (-\lambda_a,-\Lambda,-k_L,-\Delta_L),
 \label{eq:lf-parity}
\end{equation}
with convention-dependent phases stored in the adapter.  A scalar sign table
without the helicity transformation is not an implementation of parity.

\subsection{Transfer reversal}

The authoritative recoil map obeys
\begin{equation}
 \cR_{-\Delta}
 =\mathsf{Swap}_{\rm in,out}\circ\cR_{\Delta}.
 \label{eq:recoil-transfer-reversal}
\end{equation}
This identity, combined with state complex conjugation and operator adjunction,
is the numerical origin of \cref{eq:h4-hermiticity}.

\subsection{Wilson-order-zero time reversal}

H4 contains no staple rescattering.  Its operator link is the declared
zeroth-order path identity, and the state Hamiltonian is real under the
selected antiunitary convention apart from removable basis phases.  Therefore
all genuine link-odd projections satisfy
\begin{equation}
 F_{\rm link\mbox{-}odd}^{\mathrm{H4}}(x,\kT,\DT)=0.
 \label{eq:h4-link-odd-zero}
\end{equation}
This includes Sivers-, Boer--Mulders-, and gluonic-pole-type coefficients.
The zero is a controlled truncation result, not a claim that the physical
functions vanish.

\begin{requirement}[No complex-state surrogate]
\label{req:no-complex-state-surrogate}
A complex basis phase, a complex variational tensor, or a finite numerical
\(i\eps\) may not be interpreted as a naive-\(T\)-odd phase.  H4 link-odd
outputs remain zero until the typed H5 Wilson/cut engine is attached.
\end{requirement}

\subsection{Selection-rule zero ledger}

Every exact H4 zero carries a machine-readable cause selected from:
\begin{equation}
 \left\{
 \begin{array}{l}
 \texttt{SYMMETRY\_FORBIDDEN},
 \texttt{MISSING\_OAM\_BLOCK},
 \texttt{MISSING\_FOCK\_SECTOR},\\
 \texttt{WILSON\_ORDER\_ZERO},
 \texttt{KINEMATIC\_RANK\_LOSS},
 \texttt{OPERATOR\_UNAVAILABLE}
 \end{array}
 \right\}.
 \label{eq:zero-ledger}
\end{equation}
Only the first category is an exact all-model symmetry zero.  The remaining
categories are conditional on the declared approximation.

\section{Regulated reductions and matched commuting diagrams}
\label{sec:reductions}

H4 provides regulated overlap reductions.  It does not identify those
reductions automatically with renormalized QCD distributions.  Let
\begin{equation}
 \cZ_{\LF\to\QCD}^{a,J}
 \label{eq:h4-matching-placeholder}
\end{equation}
denote the future operator-basis matching map.  The architecture is
\begin{equation}
 \begin{tikzcd}[column sep=large,row sep=large]
 W_{a/N}^{\LF,r}
   \arrow[r,"\cZ_{\LF\to\QCD}"]
   \arrow[d,"\Red_{\TMD}^{\reg}"']
 &W_{a/N}^{\QCD}(\mu,\zeta)
   \arrow[d,"\Red_{\TMD}^{\QCD}"]\\
 F_{a/N}^{\LF,r}
   \arrow[r,dashed,"\cZ_{\TMD}"]
 &F_{a/N}^{\QCD}(\mu,\zeta).
 \end{tikzcd}
 \label{eq:matching-reduction-square}
\end{equation}
H4 validates the left vertical route and records the other arrows as
requirements.

\subsection{TMD limit}

The regulated TMD matrix is
\begin{equation}
 \Phi_{a/N}^{\LF,r}(x,\kT)
 =W_{a/N}^{\LF,r}(x,\kT,\DT=0).
 \label{eq:h4-tmd-limit}
\end{equation}
No information about rapidity subtraction is inferred from the existence of
this forward limit.

\subsection{GPD limit}

The regulated zero-skewness GPD matrix is
\begin{equation}
 H_{a/N}^{\LF,r}(x,0,t)
 =\int_{\Omega_{k,r}}\dd^2\kT\,
 W_{a/N}^{\LF,r}(x,\kT,\DT),
 \qquad t=-\DT^2.
 \label{eq:h4-gpd-limit}
\end{equation}
The finite basis and transverse regulator make this integral well defined.  A
physical QCD GPD still requires the corresponding renormalized operator
matching and local-limit checks.

\subsection{PDF limit}

Two regulated routes must agree:
\begin{align}
 f_{a/N}^{(1),\LF,r}(x)
 &=\int\dd^2\kT\,
 \Phi_{a/N}^{\LF,r}(x,\kT),
 \label{eq:pdf-route-one}\\
 f_{a/N}^{(2),\LF,r}(x)
 &=H_{a/N}^{\LF,r}(x,0,0).
 \label{eq:pdf-route-two}
\end{align}
Their difference
\begin{equation}
 \delta_{\PDF}^{a,r}(x)
 =f_{a/N}^{(1),\LF,r}(x)-f_{a/N}^{(2),\LF,r}(x)
 \label{eq:pdf-route-residual}
\end{equation}
is a mandatory closure residual.

\subsection{Wigner transform}

The partonic transverse Wigner distribution is
\begin{equation}
 \rho_{a/N}^{\LF,r}(x,\kT,\bD)
 =\int\frac{\dd^2\DT}{(2\pi)^2}
 e^{-i\bD\cdot\DT}
 W_{a/N}^{\LF,r}(x,\kT,\DT).
 \label{eq:h4-wigner-transform}
\end{equation}
The transform is performed componentwise on the full helicity matrix before
named projections when possible.  This makes target-spin and active-parton
interference traceable in phase space.

\subsection{Rank-aware TMD impact transform}

The independent TMD transform is
\begin{equation}
 \widetilde F^{(m)}(x,b_{\rm TMD})
 =2\pi i^m
 \int_0^\infty k_T\dd k_T\,
 J_{|m|}(b_{\rm TMD}k_T)
 N_m(k_T)F^{(m)}(x,k_T),
 \label{eq:h4-rank-bessel}
\end{equation}
with its reference mass, extracted powers, and Fourier phase stored in the
rank identity.  The Wigner coordinate \(\bD\) and TMD separation \(\bM\) are
distinct typed variables even though both have inverse-momentum dimensions.

\begin{requirement}[Reduction route closure]
\label{req:reduction-route-closure}
For every supported rank-zero species/operator channel, the TMD-to-PDF and
GPD-to-PDF routes must agree within the independently estimated quadrature,
finite-basis, and TTN errors.  A fitted route-normalization factor is
forbidden.
\end{requirement}

\section{Local-current, axial, tensor, and energy--momentum closure}
\label{sec:local-moments}

The microscopic GTMD parent is accepted only if its local moments reproduce
operators owned by the same H3 Hamiltonian identity.  This section fixes the
operator-dependent quark/antiquark signs and the gluon momentum convention.

\subsection{Vector current}

For separate positive-\(x\) quark and antiquark GPDs,
\begin{equation}
 F_1^q(t)
 =\int_0^1\dd x\,
 \left[H^q(x,0,t)-H^{\bar q}(x,0,t)\right].
 \label{eq:vector-moment}
\end{equation}
The flavor-weighted proton and neutron currents are
\begin{equation}
 F_1^N(t)
 =\sum_q e_q F_1^{q/N}(t).
 \label{eq:nucleon-vector-moment}
\end{equation}
At \(t=0\), these reproduce the H3 charge ledgers.  At nonzero transfer they
are compared with the direct H3 vector-current calculation that was withheld
from the state calibration where possible.

\subsection{Axial current}

The nonsinglet axial moment uses the charge-conjugation-even combination
\begin{equation}
 G_A^q(t)
 =\int_0^1\dd x\,
 \left[\widetilde H^q(x,0,t)+
       \widetilde H^{\bar q}(x,0,t)\right].
 \label{eq:axial-moment}
\end{equation}
It is compared to the direct axial current and to the finite-basis PCAC
closure record from H3.  The GTMD route may not alter the induced pion-pole or
pseudoscalar representation used by that record.

\subsection{Tensor current status}

For the local tensor operator, the positive-\(x\) convention gives
\begin{equation}
 g_T^q(t)
 =\int_0^1\dd x\,
 \left[H_T^q(x,0,t)-H_T^{\bar q}(x,0,t)\right]
 \label{eq:tensor-moment}
\end{equation}
under the declared charge-conjugation signature.  H4 may test this route only
if the H3 state bundle exports a Hamiltonian-consistent tensor current.  In its
absence the moment is labeled \texttt{DIRECT\_CURRENT\_COMPARATOR\_MISSING},
not normalized to a phenomenological tensor charge.

\subsection{Energy--momentum tensor}

For quarks and antiquarks,
\begin{equation}
 A_q(t)
 =\int_0^1\dd x\,x
 \left[H^q(x,0,t)+H^{\bar q}(x,0,t)\right],
 \label{eq:quark-emt-moment}
\end{equation}
while with the project convention
\begin{equation}
 H^g(x,0,0)=xg(x),
 \end{equation}
the gluon contribution is
\begin{equation}
 A_g(t)
 =\int_0^1\dd x\,H^g(x,0,t).
 \label{eq:gluon-emt-moment}
\end{equation}
Using the quark Mellin power for the gluon route is a type error.

The forward momentum closure is
\begin{equation}
 \sum_q A_q(0)+A_g(0)=1+\delta_{\rm mom}^{(r)}.
 \label{eq:emt-momentum-closure}
\end{equation}
The finite-basis residual is compared with the direct H3 longitudinal-momentum
ledger.

\subsection{Moment-closure manifest}

Every local comparison records
\begin{equation}
 \mathfrak m=
 \left(
 \begin{array}{l}
 \text{operator identity},\text{GTMD moment},
 \text{direct matrix element},\\
 \text{charge-conjugation signature},
 \text{quadrature error},\text{TTN error},\\
 \text{resolution residual},
 \text{matching status},
 \text{difference}
 \end{array}
 \right).
 \label{eq:moment-manifest}
\end{equation}
A moment is not accepted if the direct comparator uses a different Hamiltonian,
current, state member, regulator, or phase convention.

\section{Wigner phase space, orbital angular momentum, and spin--orbit closure}
\label{sec:wigner-oam}

\subsection{Canonical phase-space moment}

The transverse phase space is \(T^*\mathbb R_T^2\) with coordinates
\((\bD,\kT)\) and symplectic form
\begin{equation}
 \omega_T
 =\dd b_{\Delta x}\wedge\dd k_x
 +\dd b_{\Delta y}\wedge\dd k_y.
 \label{eq:h4-symplectic-form}
\end{equation}
The rotation moment map is
\begin{equation}
 \mu_{L_z}(\bD,\kT)
 =(\bD\times\kT)_z.
 \label{eq:h4-moment-map}
\end{equation}
For the appropriate unpolarized-parton/longitudinal-target Wigner projection,
\begin{equation}
 L_z^a
 =\int\dd x\,\dd^2\kT\,\dd^2\bD\,
 (\bD\times\kT)_z
 \rho_{LU}^a(x,\kT,\bD).
 \label{eq:wigner-oam}
\end{equation}
The path and scheme identity of the operator remain attached to the result
\cite{Lorce2011}.

\subsection{Transfer-derivative route}

Using
\begin{equation}
 b_{\Delta i}e^{-i\bD\cdot\DT}
 =i\frac{\partial}{\partial\Delta_i}
 e^{-i\bD\cdot\DT},
\end{equation}
\cref{eq:wigner-oam} is equivalent, under the declared boundary conditions,
to a transfer derivative of the GTMD matrix at \(\DT=0\):
\begin{equation}
 L_z^a
 =-i\epsilon_T^{ij}
 \int\dd x\,\dd^2\kT\,
 k_T^j
 \left.
 \frac{\partial}{\partial\Delta_T^i}
 W_{LU}^a(x,\kT,\DT)
 \right|_{\DT=0}.
 \label{eq:derivative-oam}
\end{equation}
The derivative is evaluated by analytic basis derivatives, automatic
differentiation, or symmetric finite differences with a convergence study.
A one-sided derivative at a grid boundary is not accepted without an
independent error estimate.

\subsection{GTMD-coordinate adapters}

In standard spin-\(\tfrac12\) GTMD conventions, canonical OAM and spin--orbit
correlations are associated with particular combinations commonly denoted
\(F_{1,4}\) and \(G_{1,1}\), subject to sign and normalization conventions
\cite{Kanazawa2014}.  H4 defines typed adapters
\begin{align}
 L_z^q&=\cL_F[F_{1,4}^q],\label{eq:f14-adapter}\\
 C_z^q&=\cL_G[G_{1,1}^q],
 \label{eq:g11-adapter}
\end{align}
whose mass factors, signs, Fourier phases, and Wilson-path identities are
stored explicitly.  These equations are validation identities, not license to
replace the full matrix by one scalar coefficient.

\subsection{Direct H3 comparison}

The third route is the direct regulated canonical operator in the H3 state,
\begin{equation}
 L_{z,\rm direct}^{a,r}
 =\langle\Psi_N^{(r)}|\widehat L_z^{a,r}|\Psi_N^{(r)}\rangle.
 \label{eq:direct-oam}
\end{equation}
H4 requires closure among
\begin{equation}
 L_{z,\rm Wigner}^{a,r},\qquad
 L_{z,\rm derivative}^{a,r},\qquad
 L_{z,F_{1,4}}^{a,r},\qquad
 L_{z,\rm direct}^{a,r},
 \label{eq:oam-four-routes}
\end{equation}
for every route supported by the available operator basis.

\begin{remark}[Canonical versus kinetic OAM]
The H4 straight-link/zeroth-Wilson-order result is a regulated canonical
phase-space quantity in the selected light-front representation.  Its relation
to a kinetic Ji-type decomposition or a staple-link Jaffe--Manohar quantity
requires the corresponding gauge-link, torque, and LF-to-QCD matching analysis.
H4 does not identify these definitions.
\end{remark}

\subsection{Spin--orbit correlation}

The analogous spin--orbit moment is
\begin{equation}
 C_z^a
 =\int\dd x\,\dd^2\kT\,\dd^2\bD\,
 (\bD\times\kT)_z
 \rho_{UL}^a(x,\kT,\bD),
 \label{eq:spin-orbit-wigner}
\end{equation}
where \(UL\) denotes unpolarized target and longitudinally polarized active
parton in the project convention.  Its direct comparator is the appropriately
regulated symmetrized spin--orbit operator in H3.

\section{Positivity, nonforward bounds, and quasiprobability status}
\label{sec:positivity}

\subsection{Forward Gram positivity}

At \(\DT=0\), Wilson order zero, and within the probability interpretation of
the cut overlap, the joint matrix has a Gram form
\begin{equation}
 \Phi_{\alpha\beta}(x,\kT)
 =\sum_X A_{\alpha X}(x,\kT)
 A_{\beta X}(x,\kT)^*.
 \label{eq:forward-gram}
\end{equation}
Therefore
\begin{equation}
 \Phi=\Phi^\dagger,\qquad \Phi\succeq0.
 \label{eq:forward-psd}
\end{equation}
H4 checks the minimum eigenvalue, principal minors, trace, and reconstruction
from the microscopic amplitudes.  Negative eigenvalues are not clipped.

\subsection{Off-forward overlap bound}

At \(\DT\ne0\), the matrix connects distinct momentum fibers and need not be
positive semidefinite.  For a bounded active kernel,
\begin{equation}
 \abs{W_{\beta\alpha}(\DT)}^2
 \leq
 \norm{\cK_{\beta\alpha}}_{\rm op}^2
 \norm{A_{\beta}^{\rm out}}^2
 \norm{A_{\alpha}^{\rm in}}^2.
 \label{eq:nonforward-cs-bound}
\end{equation}
The right-hand side uses the exact incoming and outgoing wave-function norms
at the shifted fibers.  Imposing forward PSD at nonzero transfer is a test
failure.

\subsection{Wigner quasidistribution}

The Wigner transform in \cref{eq:h4-wigner-transform} may be negative or
complex in off-diagonal helicity channels.  Its required conditions are the
appropriate Hermiticity, reality of physical projected moments, correct
marginals, and controlled Fourier convergence.  Pointwise positivity is not a
requirement.

\begin{requirement}[Layer-correct positivity]
\label{req:layer-correct-positivity}
The implementation must identify whether a positivity check applies to a
forward cut correlator, a reduced density matrix, an off-forward overlap, a
renormalized scheme object, or a Wigner quasidistribution.  A test from one
layer cannot be copied to another without a theorem or matching adapter.
\end{requirement}

\section{Convergence and the H4 error ledger}
\label{sec:convergence}

A small reconstruction residual at one resolution is not a microscopic
prediction.  H4 separates the following convergence axes:
\begin{equation}
 \nu_{\rm H4}=
 (K,N_{\max},b,\lambda,\cF,\chi,
 N_x,N_k,N_\Delta,N_\phi,\text{projector conditioning},\text{derivative step}).
 \label{eq:h4-convergence-axes}
\end{equation}
Here \(\chi\) is the TTN bond policy and \(\cF\) is the retained Fock content.

\subsection{Exact versus TTN closure}

For the full supported bond space,
\begin{equation}
 W_{\rm exact}^{a[J],r}
 =W_{\rm TTN,full}^{a[J],r}
 +\delta_{\rm TTN,full}^{a[J],r}.
 \label{eq:exact-ttn-gtmd}
\end{equation}
The same comparison is required for helicity matrices, scalar GTMDs, local
moments, and OAM derivatives.  For finite bond \(\chi\), errors are reported
observable by observable; energy accuracy is not a proxy for GTMD accuracy.

\subsection{Resolution and basis flow}

The comparison map from resolution \(r\) to \(r'\) acts on states and operators
before observables are compared.  The residual is
\begin{equation}
 \delta_A(r',r;z)
 =\left\|
 R_{r'\to r}
 \left[W_A^{(r')}(z)\right]
 -W_A^{(r)}(z)
 \right\|.
 \label{eq:h4-resolution-residual}
\end{equation}
Bare basis coefficients and Fock probabilities are diagnostics, not quantities
that must agree point by point across regulator schemes.

\subsection{Quadrature and projector errors}

The numerical ledger keeps separate:
\begin{align}
 \delta_{\rm quad},\quad
 \delta_{\Delta\mbox{-}grid},\quad
 \delta_{\rm Fourier},\quad
 \delta_{\rm derivative},\quad
 \delta_{\rm Gram},\quad
 \delta_{\rm reconstruction}.
 \label{eq:h4-numerical-errors}
\end{align}
A large projector condition number cannot be hidden inside the quadrature
error, and resolution variation cannot be called a statistical band.

\subsection{Observable-sensitive convergence}

The minimal release set includes convergence of:
\begin{itemize}
\item selected helicity-matrix entries;
\item all identifiable quark/antiquark coefficients;
\item each gluon polarization sector;
\item vector, axial, and EMT moments;
\item quark, antiquark, and gluon OAM ledgers;
\item forward positivity eigenvalues;
\item nonforward overlap-bound ratios;
\item TMD/GPD/PDF route residuals.
\end{itemize}
At least one low-bond or low-resolution state must demonstrate that a small
energy error can coexist with a significant GTMD, OAM, or helicity-matrix
error.

\section{Scoped replacement of the C3/C4 analytic parents}
\label{sec:replacement}

The analytic parents were essential architecture oracles.  They are not
silently deleted when the microscopic parent becomes available.  Their new
role is encoded in the provenance complex.

\subsection{Replacement relations}

The primary relations are
\begin{equation}
 \begin{array}{rcl}
 \texttt{C3\_ANALYTIC\_COMMON\_PARENT}
 &\xrightarrow{\texttt{BENCHMARKS}}&
 \texttt{H4\_MICROSCOPIC\_COMMON\_PARENT},\\[1mm]
 \texttt{C4\_ANALYTIC\_SEA\_GLUON\_PARENT}
 &\xrightarrow{\texttt{BENCHMARKS}}&
 \texttt{H4\_MICROSCOPIC\_COMMON\_PARENT},\\[1mm]
 \texttt{H4\_MICROSCOPIC\_COMMON\_PARENT}
 &\xrightarrow{\texttt{REPLACES\_WITHIN\_SCOPE}}&
 \texttt{C3/C4\_ANALYTIC\_PARENTS}.
 \end{array}
 \label{eq:h4-replacement-relations}
\end{equation}

The replacement scope is the conjunction of:
\begin{equation}
 \left\{
 \begin{array}{l}
 \xi=0,\quad\text{Wilson order}=0,\\
 a\in\{u,d,\bar u,\bar d,g\},\\
 \text{supported H3 plan/member/resolution},\\
 \text{supported operator and kinematic domain},\\
 \text{passed H4 closure and convergence gates}.
 \end{array}
 \right.
 \label{eq:h4-replacement-scope}
\end{equation}
Outside this scope, the analytic object may remain the only available oracle,
but it is never promoted to the microscopic result.

\subsection{No calibration to the analytic pilot}

Gaussian widths, sector probabilities, spin amplitudes, and synthetic phases
used by C3/C4 are not H4 calibration data.  The microscopic H3 parent may
disagree numerically.  The comparison report classifies differences as:
\begin{equation}
 \Delta W
 =\Delta W_{\rm state}
 +\Delta W_{\rm regulator}
 +\Delta W_{\rm operator}
 +\Delta W_{\rm basis}
 +\Delta W_{\rm numerical}.
 \label{eq:pilot-difference}
\end{equation}
A disagreement is not repaired with a GTMD-specific coefficient.

\subsection{Reversibility and production isolation}

The original C3/C4 nodes, manifests, tests, and source hashes remain immutable.
The H4 branch is disjoint from the accepted 216-route production root.  The
replacement edge is active only in an H4 validation compilation.  Removing H4
must restore the analytic benchmark path without changing production bytes.

\begin{requirement}[No mixed parent]
\label{req:no-mixed-parent}
An analytic parent and the H4 microscopic parent may be compared but may not be
added, averaged, spliced in kinematics, or used as mutual normalization within
one state member.
\end{requirement}

\section{Assumption compiler and provenance two-complex}
\label{sec:compiler}

\subsection{H4 assumption bundles}

An H4 \AssumptionBundle{} extends the H3 identity with:
\begin{itemize}
\item exact or TTN state representation and bond policy;
\item transfer grid and frame convention;
\item quark, antiquark, and gluon operator-basis adapters;
\item quadrature and derivative policy;
\item projector rank and conditioning thresholds;
\item Wilson order fixed to zero;
\item replacement-scope policy;
\item convergence and readiness obligations.
\end{itemize}

The primary branches inherit the H3 plans, including the chiral interaction
on/off choice.  Results from mutually exclusive branches are compared as
complete predictions and are never added.

\subsection{Prediction plan}

The H4 \PredictionPlan{} contains the directed chain
\begin{equation}
\begin{aligned}
 \StateBundle&\longrightarrow\MomentumFiber,
 \\
 \StateEvaluator&\longrightarrow\OverlapKernel,
 \\
 \HelicityMatrix&\longrightarrow\ProjectorManifest,
 \\
 \ParentBundle&\longrightarrow\MomentManifest
 \longrightarrow\WignerManifest.
\end{aligned}
\end{equation}
Every edge has typed source and target identities.  A missing readiness field
produces a structured compilation failure.

\subsection{Two-cells}

The H4 \ProvComplex{} contains finite two-cells for:
\begin{itemize}
\item exact versus full-bond TTN representation;
\item direct versus sequential PDF reduction;
\item GTMD-moment versus direct-current routes;
\item Wigner versus transfer-derivative OAM routes;
\item microscopic versus analytic parent within the declared scope;
\item explicit versus induced sector descriptions inherited from H2/H3.
\end{itemize}
A closed cycle without a declared equivalence, subtraction, or alternative
relation is an audit failure.

\section{Formal software interfaces}
\label{sec:software}

\begin{longtable}{L{0.27\textwidth}L{0.65\textwidth}}
\toprule
Object & Responsibility\\
\midrule
\endhead
\MomentumFiber & Incoming or outgoing hadron momentum fiber, helicity basis,
normalization, regulator, and phase identity.\\
\RecoilMap & Unique active/spectator zero-skewness transverse recoil map and
transfer-reversal tests.\\
\texttt{StateEvaluator} & Exact/TTN continuous momentum-space amplitudes, basis
functions, derivatives, and representation error.\\
\texttt{GTMDOverlapKernel} & Common diagonal H3 overlap with active-slot, sector, color,
helicity, recoil, and quadrature identity.\\
\texttt{HelicityMatrix} & Complete quark, antiquark, or gluon target--parton matrix
before scalar projection.\\
\texttt{CovariantBasis} & Versioned \(F\), \(G\), \(H\), or gluon tensor
basis including phases, masses, ranks, and symmetry classes.\\
\texttt{GramProjectorEngine} & Basis evaluation matrix, Gram rank,
conditioning, dual projectors, reduced basis, and reconstruction.\\
\texttt{ProjectorManifest} & Per-point projector rank, singular values, condition
number, identifiable combinations, and residuals.\\
\texttt{GTMDParentBundle} & Common microscopic quark, antiquark, and gluon matrices,
projected coefficients, matching statuses, and member identity.\\
\MomentManifest & TMD/GPD/PDF/current/EMT routes, charge-conjugation
signatures, direct-operator comparison, and residuals.\\
\texttt{WignerOAMManifest} & Transfer transform, derivative policy, OAM/spin--orbit
routes, Fourier errors, and gauge-path status.\\
\ConvergenceManifest & Basis, Fock, TTN, quadrature, transfer, projector, and
operator convergence axes.\\
\texttt{ReplacementManifest} & Scoped analytic-parent replacement and rollback
certificate.\\
\PredictionPlan & Assumption-conditioned executable graph and readiness
gates.\\
\ProvComplex & Equivalence, replacement, count-once, and alternative
relations among result paths.\\
\bottomrule
\end{longtable}

\subsection{Readiness statuses}

H4 may issue:
\begin{verbatim}
MICROSCOPIC_NONZERO_TRANSFER_PARENT_VALIDATED
T_EVEN_PROJECTOR_CLOSURE_VALIDATED
REGULATED_TMD_GPD_PDF_ROUTE_CLOSED
LOCAL_CURRENT_EMT_MOMENT_CLOSED
WIGNER_OAM_ROUTE_BENCHMARKED
ANALYTIC_PARENT_REPLACED_WITHIN_SCOPE
\end{verbatim}
It may not issue:
\begin{verbatim}
PHYSICAL_GTMD
PHYSICAL_TMD
WILSON_READY
NUCLEAR_MATCHING_READY
LF_TO_QCD_MATCHING_READY
EVOLUTION_READY
PROCESS_READY
INFERENCE_READY
\end{verbatim}

\section{Formal requirements}
\label{sec:requirements}

The implementation shall assign stable IDs to the following requirements.

\begin{longtable}{L{0.13\textwidth}L{0.79\textwidth}}
\toprule
ID & Requirement\\
\midrule
\endhead
V11.1 & Use distinct incoming and outgoing momentum fibers in the symmetric
zero-skewness frame.\\
V11.2 & Use the single authoritative active/spectator recoil map.\\
V11.3 & Prove intrinsic closure, physical transfer, spectator preservation,
unit Jacobian, forward identity, and transfer reversal.\\
V11.4 & Evaluate the H3 state in continuous momentum coordinates from exact
and TTN representations.\\
V11.5 & Preserve state, member, regulator, assumption, phase, sector, color,
permutation, helicity, and OAM identities.\\
V11.6 & Use one common overlap kernel for quarks, antiquarks, and gluons.\\
V11.7 & Select antiquarks from direct positive-\(x\) active slots.\\
V11.8 & Keep off-diagonal Fock overlaps unavailable unless a named physical
source is supplied.\\
V11.9 & Store the complete quark target--parton helicity matrix before
projection.\\
V11.10 & Store the complete direct-antiquark helicity matrix and a typed
charge-conjugation adapter.\\
V11.11 & Implement the declared four \(F\), four \(G\), and eight \(H\)
leading-twist quark GTMD covariants.\\
V11.12 & Construct projectors from the basis Gram matrix.\\
V11.13 & Report Gram rank, singular values, condition number, and null-space
combinations at every kinematic point.\\
V11.14 & Use a reduced identifiable basis at degenerate kinematics.\\
V11.15 & Prohibit pseudoinverse-based claims of fully separated coefficients
at rank-deficient points.\\
V11.16 & Store the complete microscopic gluon transverse tensor and helicity
matrix.\\
V11.17 & Retain both \(qqqg\) color outer multiplicities through every gluon
projection.\\
V11.18 & Reconstruct trace, helicity-antisymmetric, and symmetric-traceless
gluon sectors exactly.\\
V11.19 & Determine the supported gluon covariant rank from the actual Gram
matrix.\\
V11.20 & Enforce exact Wilson-order-zero link-odd null tests.\\
V11.21 & Test Hermiticity, light-front parity, transfer reversal, and T-even
reality classes on matrices and coefficients.\\
V11.22 & Expose TMD, GPD, and PDF reductions as views of one parent.\\
V11.23 & Close direct and sequential PDF routes.\\
V11.24 & Close vector-current moments with \(q-\bar q\).\\
V11.25 & Close axial moments with \(q+\bar q\).\\
V11.26 & Close tensor moments with the declared charge-conjugation signature.\\
V11.27 & Close quark and gluon EMT moments with their distinct Mellin
conventions.\\
V11.28 & Compare every moment route with a Hamiltonian-consistent direct H3
operator.\\
V11.29 & Keep regulated reductions distinct from QCD-matched operators.\\
V11.30 & Keep \(\bD\) and \(\bM\) as distinct coordinate types.\\
V11.31 & Implement the partonic Wigner transform with an explicit finite-range
and interpolation error.\\
V11.32 & Compare Wigner, transfer-derivative, and direct-operator OAM routes.\\
V11.33 & Treat \(F_{1,4}\) and \(G_{1,1}\) OAM/spin--orbit relations as
versioned convention adapters.\\
V11.34 & Enforce forward Gram positivity without clipping.\\
V11.35 & Use a Cauchy--Schwarz overlap bound, not PSD, at nonzero transfer.\\
V11.36 & Never impose pointwise positivity on Wigner functions.\\
V11.37 & Report independent basis, Fock, TTN, quadrature, transfer, derivative,
projector, current, regulator, and matching errors.\\
V11.38 & Require observable-sensitive TTN bond convergence.\\
V11.39 & Replace C3/C4 analytic parents only within a typed, reversible scope.\\
V11.40 & Never tune the microscopic parent to an analytic pilot curve.\\
V11.41 & Preserve the assumption-conditioned plan and provenance two-complex.\\
V11.42 & Keep production, nuclear, matching, evolution, process, and inference
roots fail-closed.\\
\bottomrule
\end{longtable}

\section{Benchmark and falsification program}
\label{sec:benchmarks}

\subsection{H4-A: recoil and momentum fibers}

Evaluate all H3 sectors at several nonzero transfers and verify exact
intrinsic closure, spectator preservation, unit Jacobian, transfer reversal,
and the forward identity.  Inject wrong active and spectator signs, a local
species-specific recoil formula, and an incoming/outgoing fiber swap.

\subsection{H4-B: exact and TTN momentum amplitudes}

Compare exact basis sums and full-bond TTN amplitudes pointwise and in norm.
Test phases, recoupling, color multiplicity, fermion signs, and derivatives.
A low-bond member must show an observable-sensitive amplitude error even when
its energy is accurate.

\subsection{H4-C: scalar overlap oracle}

Recover the C3 one- and two-body analytic limits by inserting the corresponding
controlled states into the H4 evaluator.  This is an algebraic oracle, not a
fit to the H3 state.

\subsection{H4-D: quark sixteen-function closure}

At generic \((\kT,\DT)\), generate a known linear combination of the sixteen
covariants, recover every coefficient with the Gram projectors, and reconstruct
the full helicity matrix.  Repeat with the microscopic H3 matrix and report
conditioning.

\subsection{H4-E: degeneracy and rank loss}

Approach \(\DT=0\), \(\kT=0\), and collinear orientations.  Verify the
expected Gram-rank reduction and the explicit null-space combinations.  An
injected pseudoinverse-based full-coefficient claim must fail.

\subsection{H4-F: direct antiquark closure}

Activate \(\bar u\) and \(\bar d\) slots independently.  Verify positive
support, pair-sector ancestry, anti-fundamental color action, direct helicity
matrix reconstruction, and operator-dependent moment signs.  Negative-\(x\)
copying must fail.

\subsection{H4-G: gluon tensor and helicity closure}

From one microscopic gluon parent, reconstruct trace,
helicity-antisymmetric, and symmetric-traceless components in both color
outer-multiplicity ancestries.  Test the full tensor and helicity conversions.

\subsection{H4-H: discrete symmetries}

Test Hermiticity at \(\pm\DT\), light-front parity, transfer reversal,
reality classes, and exact T-odd nulls at Wilson order zero.  An array-only
future/past label swap must not count as an antiunitary test.

\subsection{H4-I: marginal closure}

Compare GTMD-to-TMD, GTMD-to-GPD, direct PDF, and GPD-forward PDF routes for
all supported species and several transfers.  Keep quadrature and projector
residuals separate.

\subsection{H4-J: local moments}

Compare vector, axial, tensor, and EMT moments with direct H3 operators.
Inject wrong quark--antiquark signs and the wrong gluon Mellin convention.

\subsection{H4-K: Wigner and OAM closure}

Compare the Wigner phase-space moment, transfer derivative, direct H3 OAM,
and supported \(F_{1,4}\)/\(G_{1,1}\) adapters.  Vary transfer range,
derivative method, and Fourier grid.

\subsection{H4-L: positivity and overlap bounds}

Test forward Gram positivity, nonforward Cauchy--Schwarz bounds, and permitted
Wigner negativity.  Inject forward clipping, nonforward PSD enforcement, and
pointwise Wigner positivity.

\subsection{H4-M: multi-axis convergence}

Run the parent and selected moments across at least three basis resolutions,
multiple TTN bonds, and refined quadrature/transfer grids.  Demonstrate that
convergence in energy does not substitute for GTMD or OAM convergence.

\subsection{H4-N: analytic-parent replacement}

Compile the scoped H4 replacement, verify that C3/C4 remain immutable, and
show that unsupported species, operators, transfers, or failed gates do not
silently inherit an analytic result.

\subsection{H4-O: assumption branches}

Execute the supported H3 plans as complete alternatives.  Compare their
microscopic parents and reductions without adding them.  A branch-specific
operator or state cannot be consumed by another branch.

\section{Negative-test catalogue}
\label{sec:negative}

The implementation shall provide stable IDs for at least the following one
hundred classes; grouping does not limit the final count.

\begin{longtable}{L{0.18\textwidth}L{0.72\textwidth}}
\toprule
Group & Required injected failures\\
\midrule
\endhead
Fibers and recoil & Wrong incoming/outgoing momentum, nonzero skewness,
transfer sign error, active factor error, spectator factor error, local recoil
reimplementation, broken intrinsic closure, nonunit Jacobian, spectator
momentum change, transfer reversal failure.\\
State evaluation & Wrong basis normalization, oscillator scale called TMD
width, exact/TTN phase mismatch, lost color multiplicity, lost fermion sign,
wrong center-of-mass factor, unsupported interpolation, derivative without
phase alignment, full-bond mismatch.\\
Overlap & Wrong active species, duplicate active slot, missing sector ancestry,
off-diagonal sector without source, mismatched regulator, state member mixed,
quark/gluon normalization copied, missing spectator match, Wilson order
nonzero in the H4 central route.\\
Quark basis & Missing or duplicated \(F\), \(G\), or \(H\) covariant, wrong
mass or \(P^+\) factor, wrong \(i\) phase, wrong transverse epsilon,
misordered helicity basis, convention adapter absent, reconstruction failure.\\
Projectors & Wrong Gram metric, hidden rank loss, singular inverse accepted,
pseudoinverse promoted to full identifiability, unreported null space, unstable
condition number ignored, projector from another kinematic point reused,
coefficient accepted without matrix reconstruction.\\
Antiquarks & Negative-\(x\) copy, fundamental rather than anti-fundamental
color action, wrong charge-conjugation signature, lost pair flavor, unsupported
antiquark sector, quark moment sign copied to axial or tensor operator.\\
Gluons & One color outer multiplicity dropped, trace normalization wrong,
helicity-antisymmetric sign wrong, symmetric-traceless trace nonzero, gluon
helicity conversion wrong, fixed scalar count assumed despite Gram rank,
link-odd coefficient nonzero at Wilson order zero.\\
Symmetries & Hermiticity failure, parity adapter omitted, transfer reversal
implemented as array reversal only, T-even reality class broken, future/past
labels used as a physical phase at order zero.\\
Marginals & TMD/GPD/PDF routes use different parents, direct and sequential
PDF mismatch hidden, collinear integral applied to wrong rank, current moment
normalized independently, regulated result labeled matched.\\
Moments & Vector \(q+\bar q\), axial \(q-\bar q\), tensor wrong sign, quark
Mellin power used for gluon, direct operator from another Hamiltonian, PCAC or
Ward residual repaired with a GTMD coefficient.\\
Wigner/OAM & \(\bD\) confused with \(\bM\), Fourier sign wrong, finite range
unreported, transfer derivative one-sided without study, OAM paths mixed,
\(F_{1,4}\) sign hard-coded, Wigner negativity clipped.\\
Positivity & Forward negative eigenvalue clipped, nonforward matrix required
PSD, overlap kernel norm omitted, Wigner pointwise positivity imposed,
finite-order matched positivity assumed without theorem.\\
Convergence & Energy-only TTN acceptance, basis and quadrature errors merged,
state vectors compared across regulators without map, required OAM sector
truncated silently, derivative step fixed without study, projector error hidden.\\
Replacement and provenance & Analytic parent overwritten, replacement scope absent,
unsupported domain falls back to analytic number, analytic width fitted to
microscopic state, branch alternatives added, equivalence cycle missing,
rollback disabled.\\
Downstream gates & Wilson, nuclear, LF-to-QCD, evolution, process, inference,
or production promotion; physical GTMD/TMD label; accepted registry mutation;
authoritative artifact or prior oracle changed.\\
\bottomrule
\end{longtable}

\section{Staged H4 implementation program}
\label{sec:stages}

\subsection{H4.0: fibers and continuous state evaluator}

Implement momentum fibers, reuse the recoil authority, evaluate exact and TTN
H3 states at shifted coordinates, and pass H4-A--C before constructing any
full GTMD basis.

\subsection{H4.1: common microscopic overlap}

Implement the common diagonal overlap kernel and active quark, antiquark, and
gluon selectors.  Close normalization, support, color, permutation, and
forward limits.

\subsection{H4.2: helicity matrices and projectors}

Build the quark/antiquark sixteen-covariant basis, gluon tensor/helicity basis,
Gram projector engine, degeneracy handling, and symmetry tests.

\subsection{H4.3: reductions and moments}

Implement TMD/GPD/PDF views, current and EMT moments, Wigner transform, transfer
derivatives, OAM/spin--orbit adapters, positivity, and nonforward bounds.

\subsection{H4.4: convergence and scoped replacement}

Run multi-resolution, multi-bond, and numerical convergence studies.  Activate
the microscopic replacement of C3/C4 only after all declared gates pass.

\subsection{H4.5: release and handoff}

Generate the H4 state-parent bundle, projector and moment manifests,
convergence report, replacement certificate, regression evidence, and the
microscopic-parent handoff required by the future H5 Wilson package.

\section{Risks and safeguards}
\label{sec:risks}

\begin{longtable}{L{0.24\textwidth}L{0.32\textwidth}L{0.36\textwidth}}
\toprule
Risk & Failure mode & Required safeguard\\
\midrule
\endhead
False off-forward dynamics & Recoil or state phases are inserted locally and
produce convention-dependent structure. & One typed recoil authority,
momentum fibers, transfer reversal, and phase-aligned exact/TTN tests.\\
Projector overclaim & Rank-deficient kinematics are inverted and all GTMDs are
reported separately. & Gram-rank and null-space manifest; reduced basis;
fail-closed identifiability.\\
Tensor-network blindness & Low bond reproduces energy but erases the only OAM
or helicity interference supporting a function. & Observable-sensitive bond
convergence and required symmetry-sector retention.\\
Analytic-parent anchoring & The microscopic state is tuned to reproduce the
Gaussian pilot. & Pilots are algebraic oracles only; no pilot likelihood or
GTMD-specific correction.\\
Positivity misuse & PSD is imposed off forward or Wigner negativity is
clipped. & Forward Gram theorem, nonforward overlap bounds, Wigner
quasiprobability status.\\
Moment normalization drift & Each form factor or moment is rescaled
independently. & Hamiltonian-consistent direct operators and common state
normalization.\\
Gluon incompleteness & One color multiplicity or polarization sector is lost
before projection. & Microscopic ancestry retained and tensor/helicity
reconstruction tested.\\
Matching overreach & Regulated overlaps are called physical GTMDs or TMDs. &
Explicit unresolved UV, rapidity, soft, link-shortening, and evolution status.\\
Premature Wilson phase & Future/past labels create a phase at Wilson order
zero. & Exact link-odd null tests and inactive C5/C6 handoff.\\
Numerical convergence laundering & One converged axis hides another. &
Separate basis, Fock, TTN, quadrature, derivative, projector, and operator
residuals.\\
\bottomrule
\end{longtable}

\section{Formal acceptance criteria}
\label{sec:acceptance}

H4 is complete for its declared scope only if:
\begin{enumerate}
\item the exact C10/H3 baseline and all prior regression gates reproduce;
\item incoming/outgoing fibers and the recoil map pass all exact identities;
\item exact and full-bond TTN momentum-space states agree;
\item one overlap kernel supplies all supported species;
\item direct positive-\(x\) antiquarks remain distinct from quarks;
\item complete quark and antiquark helicity matrices are retained;
\item the generic quark sixteen-function basis reconstructs exactly;
\item rank loss and null combinations are reported at degenerate kinematics;
\item the gluon tensor and helicity representations reconstruct exactly;
\item both microscopic gluon color multiplicities remain traceable;
\item Hermiticity, parity, transfer reversal, and Wilson-order-zero nulls pass;
\item TMD, GPD, PDF, current, and EMT reductions commute to tolerance;
\item vector, axial, tensor, and gluon moment conventions are correct;
\item Wigner, derivative, direct-OAM, and supported GTMD-adapter routes are
consistent within separated errors;
\item forward positivity and nonforward overlap bounds pass without clipping;
\item basis, Fock, TTN, quadrature, transfer, derivative, and projector
convergence are separately demonstrated;
\item C3/C4 remain immutable and are replaced only within a certified scope;
\item no microscopic coefficient is fitted to an analytic parent output;
\item every result has a content-addressed assumption and member identity;
\item all new negative injections are detected with stable diagnostics;
\item production registry, provenance, composition, and authoritative
artifacts remain byte-identical;
\item Wilson, nuclear, matching, evolution, process, and inference roots remain
fail-closed;
\item all schemas, manifests, source hashes, and handoff notes are complete.
\end{enumerate}

\section{Recommended decision and H5 handoff}
\label{sec:decision}

The adopted H4 object is a
\begin{quote}
\emph{regulated, zero-skewness, Wilson-order-zero microscopic quark,
antiquark, and gluon GTMD helicity-matrix bundle generated from one H3
Hamiltonian eigenstate, with Gram-certified T-even projections, commuting
marginals and local moments, Wigner/OAM closure, observable-sensitive TTN
convergence, and a scoped replacement of analytic overlap pilots.}
\end{quote}

A successful H4 package supplies the required input for:
\begin{quote}
\textbf{H5: microscopic Wilson dynamics on the H4 parent,}
including physical spectral support, quark and antiquark link-odd helicity
matrices, distinct Sivers and Boer--Mulders projections, independent active-
gluon \(f/d\) channels, soft/rapidity accounting, and Wilson-order/Fock-order
compatibility.
\end{quote}
H5 must not begin by assigning phases to the scalar H4 coefficients.  It must
act on the common matrix parent and use the same microscopic state, gauge
couplings, color multiplicities, and OAM blocks.

\appendix

\section{Compact mathematical dictionary}

\begin{tabularx}{\textwidth}{L{0.30\textwidth}Y}
\toprule
Physical concept & Mathematical/software object\\
\midrule
Incoming/outgoing hadron states & Separate \MomentumFiber{} objects\\
Zero-skewness transfer & Versioned \RecoilMap{}\\
Exact or TTN wave function & \StateEvaluator{}\\
Common partonic matrix element & \OverlapKernel{}\\
Target--parton spin information & \HelicityMatrix{}\\
Named GTMD basis & \texttt{CovariantBasis}\\
Identifiability & Gram rank and null space\\
Named coefficient extraction & Dual Gram projectors\\
Quark--antiquark signs & Operator charge-conjugation signature\\
Gluon polarization & Trace, antisymmetric, symmetric-traceless tensors\\
TMD/GPD/PDF/current views & Typed \(\Red\) maps\\
Partonic phase space & \(\rho_W(x,\kT,\bD)\)\\
OAM & Wigner moment/transfer derivative/direct operator\\
Forward positivity & Gram positive-semidefinite matrix\\
Nonforward constraint & Cauchy--Schwarz overlap bound\\
Analytic-to-microscopic transition & \ReplacementManifest{}\\
Conditional calculation & \AssumptionBundle{} and \PredictionPlan{}\\
Composition audit & \ProvComplex{}\\
\bottomrule
\end{tabularx}

\section{Microscopic parent schema}

\begin{verbatim}
{
  "parent_bundle_id": "...",
  "state_bundle_id": "...",
  "assumption_bundle_id": "...",
  "resolution_id": "...",
  "member_id": "...",
  "wilson_order": 0,
  "skewness": 0.0,
  "incoming_fiber_id": "...",
  "outgoing_fiber_id": "...",
  "recoil_map_id": "...",
  "state_representation": "EXACT | TTN",
  "bond_dimension": "...",
  "species": ["u", "d", "ubar", "dbar", "g"],
  "operator_ids": ["..."],
  "helicity_matrix_assets": {"...": "..."},
  "projector_manifest_ids": ["..."],
  "marginal_assets": {"TMD": "...", "GPD": "...", "PDF": "..."},
  "moment_manifest_id": "...",
  "wigner_oam_manifest_id": "...",
  "convergence_manifest_id": "...",
  "matching_status": [
    "LINK_SHORTENING_REQUIRED",
    "UV_MATCHING_REQUIRED",
    "RAPIDITY_SOFT_MATCHING_REQUIRED",
    "NO_EVOLUTION_APPLIED",
    "NO_PROCESS_MAP_APPLIED"
  ],
  "readiness": "MICROSCOPIC_NONZERO_TRANSFER_PARENT_VALIDATED"
}
\end{verbatim}

\section{Projector manifest schema}

\begin{verbatim}
{
  "projector_manifest_id": "...",
  "basis_id": "QUARK_FGH_V1 | GLUON_COVARIANT_V1",
  "species": "...",
  "kinematics": {"x": 0.0, "kT": [0.0, 0.0], "DeltaT": [0.0, 0.0]},
  "helicity_basis_id": "...",
  "gram_rank": 0,
  "basis_dimension": 0,
  "singular_values": [0.0],
  "condition_number": 0.0,
  "rank_tolerance": 0.0,
  "identifiable_coefficients": ["..."],
  "null_combinations": ["..."],
  "reconstruction_residual": 0.0,
  "coefficient_covariance_numerical": [[0.0]],
  "status": "FULL_RANK | REDUCED_BASIS | UNAVAILABLE"
}
\end{verbatim}

\section{Moment and Wigner/OAM schemas}

\begin{verbatim}
{
  "moment_manifest_id": "...",
  "parent_bundle_id": "...",
  "operator_id": "VECTOR | AXIAL | TENSOR | EMT_QUARK | EMT_GLUON",
  "charge_conjugation_signature": "MINUS | PLUS | NOT_APPLICABLE",
  "direct_operator_result": "...",
  "gtmd_gpd_moment_result": "...",
  "direct_integral_result": "...",
  "quadrature_residual": 0.0,
  "projector_residual": 0.0,
  "state_representation_residual": 0.0,
  "matching_status": "REGULATED_ANALYTICALLY_UNMATCHED"
}
\end{verbatim}

\begin{verbatim}
{
  "wigner_oam_manifest_id": "...",
  "parent_bundle_id": "...",
  "species": "...",
  "bDelta_grid_id": "...",
  "fourier_convention_id": "...",
  "phase_space_result": "...",
  "transfer_derivative_result": "...",
  "direct_oam_result": "...",
  "F14_adapter_result": "...",
  "G11_adapter_result": "...",
  "transfer_range_residual": 0.0,
  "derivative_step_residual": 0.0,
  "fourier_residual": 0.0,
  "gauge_path_status": "STRAIGHT_OR_ORDER_ZERO_REGULATED"
}
\end{verbatim}

\section{Replacement and convergence schemas}

\begin{verbatim}
{
  "replacement_manifest_id": "...",
  "microscopic_parent_id": "...",
  "analytic_parent_ids": ["C3:...", "C4:..."],
  "scope": {
    "species": ["..."],
    "operators": ["..."],
    "xi": 0.0,
    "wilson_order": 0,
    "kinematic_domain_id": "...",
    "assumption_plan_id": "..."
  },
  "required_gate_ids": ["..."],
  "unsupported_scope": ["..."],
  "rollback_rule": "FAIL_CLOSED_UNAVAILABLE",
  "analytic_oracles_immutable": true
}
\end{verbatim}

\begin{verbatim}
{
  "convergence_manifest_id": "...",
  "parent_bundle_id": "...",
  "basis_tower": ["..."],
  "fock_content": ["QQQ", "QQQG", "QQQUUBAR", "QQQDDBAR"],
  "bond_dimensions": [1, 2, 4],
  "quadrature_levels": ["..."],
  "transfer_grids": ["..."],
  "derivative_steps": [0.0],
  "projector_tolerances": [0.0],
  "observable_residuals": {"...": 0.0},
  "accepted_domains": {"...": "..."},
  "unresolved_axes": ["UV_MATCHING", "RAPIDITY_SOFT", "CONTINUUM"]
}
\end{verbatim}

\small
\begin{thebibliography}{99}
\raggedright

\bibitem{Dirac1949}
P.~A.~M.~Dirac,
``Forms of relativistic dynamics,''
Rev. Mod. Phys. \textbf{21}, 392 (1949).

\bibitem{BrodskyPauliPinsky1998}
S.~J.~Brodsky, H.-C.~Pauli, and S.~S.~Pinsky,
``Quantum chromodynamics and other field theories on the light cone,''
Phys. Rept. \textbf{301}, 299 (1998),
arXiv:hep-ph/9705477.

\bibitem{Vary2010}
J.~P.~Vary et al.,
``Hamiltonian light-front field theory in a basis function approach,''
Phys. Rev. C \textbf{81}, 035205 (2010),
arXiv:0905.1411.

\bibitem{Mondal2025}
C.~Mondal et al.,
``Basis light-front quantization: Advancing a first principles approach for
the nucleon,''
arXiv:2503.21372 (2025).

\bibitem{BrodskyDiehlHwang2001}
S.~J.~Brodsky, M.~Diehl, and D.-S.~Hwang,
``Light-cone wavefunction representation of deeply virtual Compton
scattering,''
Nucl. Phys. B \textbf{596}, 99 (2001),
arXiv:hep-ph/0009254.

\bibitem{Diehl2003}
M.~Diehl,
``Generalized parton distributions,''
Phys. Rept. \textbf{388}, 41 (2003),
arXiv:hep-ph/0307382.

\bibitem{Meissner2009}
S.~Mei{\ss}ner, A.~Metz, and M.~Schlegel,
``Generalized parton correlation functions for a spin-1/2 hadron,''
JHEP \textbf{08}, 056 (2009),
arXiv:0906.5323.

\bibitem{PasquiniLorce2013}
B.~Pasquini and C.~Lorc\'e,
``The partonic structure of the nucleon from generalized transverse
momentum-dependent parton distributions,''
Phys. Rev. D \textbf{88}, 034020 (2013),
arXiv:1304.1479.

\bibitem{LorcePasquini2013}
C.~Lorc\'e and B.~Pasquini,
``Structure analysis of the generalized correlator of quark and gluon for a
spin-1/2 target,''
JHEP \textbf{09}, 138 (2013),
arXiv:1307.4497.

\bibitem{Lorce2011}
C.~Lorc\'e and B.~Pasquini,
``Quark Wigner distributions and orbital angular momentum,''
Phys. Rev. D \textbf{84}, 014015 (2011),
arXiv:1106.0139.

\bibitem{LorcePasquini2016}
C.~Lorc\'e and B.~Pasquini,
``Multipole decomposition of the nucleon transverse phase space,''
Phys. Rev. D \textbf{93}, 034040 (2016),
arXiv:1512.06744.

\bibitem{Kanazawa2014}
K.~Kanazawa, C.~Lorc\'e, A.~Metz, B.~Pasquini, and M.~Schlegel,
``Twist-2 generalized transverse-momentum dependent parton distributions and
the spin/orbital structure of the nucleon,''
Phys. Rev. D \textbf{90}, 014028 (2014),
arXiv:1403.5226.

\bibitem{Echevarria2016}
M.~G.~Echevarria, A.~Idilbi, K.~Kanazawa, C.~Lorc\'e, A.~Metz,
B.~Pasquini, and M.~Schlegel,
``Proper definition and evolution of generalized transverse momentum
distributions,''
Phys. Lett. B \textbf{759}, 336 (2016),
arXiv:1602.06953.

\bibitem{Collins2011}
J.~Collins,
\emph{Foundations of Perturbative QCD},
Cambridge University Press (2011).

\bibitem{Pobylitsa2002}
P.~V.~Pobylitsa,
``Positivity bounds on generalized parton distributions in impact parameter
representation,''
Phys. Rev. D \textbf{66}, 094002 (2002),
arXiv:hep-ph/0204337.

\bibitem{Cotogno2017}
S.~Cotogno, T.~van Daal, and P.~J.~Mulders,
``Positivity bounds on gluon TMDs for hadrons of spin \(\leq1\),''
JHEP \textbf{11}, 185 (2017),
arXiv:1709.07827.

\bibitem{Singh2010}
S.~Singh, R.~N.~C.~Pfeifer, and G.~Vidal,
``Tensor network decompositions in the presence of a global symmetry,''
Phys. Rev. A \textbf{82}, 050301(R) (2010),
arXiv:0907.2994.

\bibitem{Corboz2010}
P.~Corboz, G.~Evenbly, F.~Verstraete, and G.~Vidal,
``Simulation of interacting fermions with entanglement renormalization,''
Phys. Rev. A \textbf{81}, 010303(R) (2010),
arXiv:0904.4151.

\bibitem{Feshbach1958}
H.~Feshbach,
``Unified theory of nuclear reactions,''
Ann. Phys. \textbf{5}, 357 (1958),
and Ann. Phys. \textbf{19}, 287 (1962).

\bibitem{Marinho2008}
J.~A.~O.~Marinho, T.~Frederico, E.~Pace, G.~Salm\`e, and P.~U.~Sauer,
``Light-front Ward--Takahashi identity for two-fermion systems,''
Phys. Rev. D \textbf{77}, 116010 (2008),
arXiv:0805.0707.

\end{thebibliography}

\end{document}
